\section{NestSection Results}
\label{sec:results}

\subsection{Spine Construction Success Rate}
\label{sec:results:success}

NestSection successfully produces nested bicameral maps for all 42 compatible
states. Success is defined as: (1) all senate districts are exact unions of
house districts; (2) all house districts satisfy population balance
$\leq 0.5\%$; (3) all senate districts satisfy population balance $\leq 0.5\%$;
and (4) all districts are contiguous.

One compatibility edge case requires attention: states with small spine counts
($g \leq 4$) and large per-spine district counts ($H/g \geq 30$) can have
difficulty achieving simultaneous population balance for both chambers within
each spine super-region. This occurs because the spine super-region's population
must divide evenly into both $H/g$ and $S/g$ parts---a simultaneous balancing
constraint that is more demanding than either chamber's constraint alone.
In practice, all 42 compatible states succeed because METIS's balance tolerance
is applied at the sub-region level, not the full-state level.

\subsection{Example Maps}
\label{sec:results:examples}

Table~\ref{tab:nesting-examples} shows NestSection results for selected states.

\begin{table}[h]
\centering\small
\begin{tabular}{llrrrrrr}
\toprule
State & Ratio & $g$ & PP (house) & PP (senate) & Max dev (H) & Max dev (S) \\
\midrule
Washington & 2:1 & 49 & 0.412 & 0.389 & 0.31\% & 0.38\% \\
Minnesota  & 2:1 & 67 & 0.398 & 0.371 & 0.29\% & 0.35\% \\
Florida    & 3:1 & 40 & 0.381 & 0.361 & 0.42\% & 0.47\% \\
Ohio       & 3:1 & 33 & 0.374 & 0.349 & 0.44\% & 0.48\% \\
Vermont    & 5:1 & 30 & 0.402 & 0.377 & 0.39\% & 0.45\% \\
Massachusetts & 4:1 & 40 & 0.389 & 0.361 & 0.40\% & 0.46\% \\
New Hampshire & 50:3 & 8 & 0.388 & 0.341 & 0.48\% & 0.49\% \\
\bottomrule
\end{tabular}
\caption{NestSection results for selected states. PP = mean Polsby-Popper.
  Max dev = maximum population deviation. All states achieve $\leq 0.5\%$
  maximum deviation in both chambers. Senate PP is consistently lower than
  house PP: larger senate districts are less compact on average.}
\label{tab:nesting-examples}
\end{table}

\paragraph{Observation: senate PP $<$ house PP.}
In all 42 compatible states, senate districts have lower mean PP than house
districts drawn from the same spine. This is a structural consequence of
nesting: senate districts are composed of multiple house districts, and the
union of two or more compact convex regions is generally less compact than
either region alone. Specifically, the boundary between two house districts
that are joined to form a senate district becomes an interior boundary
of the senate district, increasing the perimeter relative to the area.
The PP penalty for nesting is approximately $0.02$--$0.04$ PP points
(senate PP = house PP $- 0.028$ on average across the 42 states).
